\subsection{Task 3.8. Kế hoạch chi tiết lập trong quá trình thực hiện}
\label{sec:t3-lower}

Theo Step Wise, kế hoạch chi tiết chỉ được lập khi có đủ thông tin. Ba kế hoạch dưới đây được lập trong giai đoạn thực hiện, bổ sung cho backlog I1 (Bảng \ref{tab:i1-lower}). Chúng chia nhỏ các hoạt động đã có trong baseline và \textbf{không làm thay đổi tổng giờ công đã duyệt}.

\subsubsection{Backlog Increment 2 (lập ngày 12/11)}

\begin{table}[H]
\centering\footnotesize
\caption{Kế hoạch chi tiết I2}\label{tab:t3-lower-i2}
\begin{tabularx}{\linewidth}{@{}p{.09\linewidth}Xp{.32\linewidth}@{}}
\toprule
\textbf{Gói} & \textbf{Phần việc chi tiết} & \textbf{Tiêu chí hoàn thành}\\\midrule
A07.1 & Mô hình một dòng đăng ký hiện hành/cặp sự kiện--sinh viên; bảng lịch sử; trigger reset điểm danh khi đăng ký lại. & Đăng ký lại giữ \texttt{id}; đúng 1 dòng lịch sử REREGISTER.\\
A07.1 & Khóa sự kiện trong transaction, kiểm tra khoa AND năm, khung giờ. & Test 6 người/1 chỗ đạt.\\
A07.2 & Đôn FIFO theo (\texttt{waitlisted\_at}, \texttt{id}); bỏ qua người không hợp lệ (\texttt{SKIP}); đôn khi tăng sức chứa. & Thứ tự đúng; người bị bỏ qua giữ vị trí.\\
A07.3 & Bảng \texttt{email\_outbox}; worker gửi sau commit; retry có giới hạn; chống gửi thư lượt cũ. & Lỗi ghi outbox làm rollback; 503 retry, 400 dừng.\\
A08 & Frontend I2: panel đăng ký, sự kiện của tôi, danh sách người tham gia. & Regression I1 và test I2 đạt.\\
\bottomrule
\end{tabularx}
\end{table}

\subsubsection{Backlog Increment 3 (lập ngày 18/11)}

\begin{table}[H]
\centering\footnotesize
\caption{Kế hoạch chi tiết I3}\label{tab:t3-lower-i3}
\begin{tabularx}{\linewidth}{@{}p{.09\linewidth}Xp{.32\linewidth}@{}}
\toprule
\textbf{Gói} & \textbf{Phần việc chi tiết} & \textbf{Tiêu chí hoàn thành}\\\midrule
A09.1 & Phiên QR 1--30 phút, tối đa một phiên mở; token 32 byte lưu SHA-256 + AES-GCM; QR SVG. & Mã cũ vô hiệu sau khi đóng hoặc hết hạn.\\
A09.1 & Tách API preview (chỉ đọc) và check-in (có version); guard ở database tại thời điểm ghi. & Xác nhận cũ trả \texttt{CONFIRMATION\_STALE}; không ghi một phần.\\
A09.1 & Trang \texttt{/scan}: camera, giải mã QR, màn xác nhận. & Không tự gửi check-in khi vừa đọc mã.\\
A09.2 & API chỉnh thủ công dùng chung, lý do bắt buộc, lịch sử snapshot trước/sau. & Báo cáo cập nhật sau khi chỉnh.\\
A10.1 & Báo cáo một snapshot; vắng chỉ tính sau \texttt{ends\_at}; mẫu số là \texttt{REGISTERED}. & Đối chiếu đúng với dữ liệu biết trước kết quả.\\
A10.2 & CSV UTF-8 BOM, chống công thức, giới hạn 10.000 dòng. & 10.001 dòng trả 413.\\
\bottomrule
\end{tabularx}
\end{table}

\subsubsection{Kế hoạch triển khai (lập ngày 25/11)}

\begin{table}[H]
\centering\footnotesize
\caption{Kế hoạch chi tiết A11}\label{tab:t3-lower-deploy}
\begin{tabularx}{\linewidth}{@{}p{.05\linewidth}Xp{.32\linewidth}@{}}
\toprule
\textbf{Bước} & \textbf{Công việc} & \textbf{Tiêu chí kiểm tra}\\\midrule
1 & Rà soát, sửa lỗi và merge toàn bộ I1--I3 vào \texttt{main}. & CI và bộ test đạt trên \texttt{main}.\\
2 & Tạo database Neon; Flyway tự tạo schema khi backend chạy. & 15 bảng nghiệp vụ; không bật dữ liệu seed.\\
3 & Tạo dịch vụ Render từ \texttt{render.yaml}; đặt secret (DB, khóa QR, Brevo). & \texttt{/api/live} và \texttt{/api/health} trả 200.\\
4 & Tạo Admin đầu tiên bằng bootstrap rồi tắt bootstrap. & Đăng nhập được; biến bootstrap đã gỡ.\\
5 & Tạo dự án Vercel với proxy \texttt{/api} sang Render. & Cookie, CSRF, đăng xuất hoạt động qua HTTPS.\\
6 & Chạy luồng end-to-end và quét QR bằng điện thoại thật. & \TBD{kết quả}.\\
\bottomrule
\end{tabularx}
\end{table}

\subsubsection{Vì sao các chi tiết này không thể lập từ đầu}

\begin{itemize}
    \item \textbf{Phụ thuộc các quyết định nghiệp vụ còn mở.} Quy tắc đôn chờ sau giờ đóng đăng ký và các mốc gửi email được nhóm chốt ngày 12/11; giờ mở QR và bước xác nhận được chốt ngày 18/11. Các quyết định này dựa trên đề xuất A3--A5 và chưa được khách hàng xác nhận. Trước các ngày này, nhóm chỉ có thể lập kế hoạch ở mức hoạt động A07, A09.
    \item \textbf{Phụ thuộc kết quả increment trước.} Thiết kế transaction của I2 dựa trên kiến trúc phân tầng hình thành sau I1; thiết kế điểm danh I3 dựa trên mô hình đăng ký của I2.
    \item \textbf{Phụ thuộc kết quả thử nghiệm.} Kế hoạch triển khai chỉ cụ thể được sau khi chọn nền tảng (CR-04) và gặp các giới hạn thực tế như Maven 429, proxy cookie.
\end{itemize}
